\documentclass[10pt,a4paper]{article}

% Shared presentation layer for the standalone R0--R7 development documents.
% Method equations remain authoritative in the canonical idea sketch.
\usepackage[a4paper,margin=18mm,headheight=14pt]{geometry}
\usepackage{fontspec}
\usepackage{xeCJK}
\xeCJKsetup{CJKspace=true}
\setmainfont{DejaVu Serif}
\setsansfont{DejaVu Sans}
\setmonofont{DejaVu Sans Mono}
\setCJKmainfont{Noto Serif CJK KR}
\setCJKsansfont{Noto Sans CJK KR}
\setCJKmonofont{Noto Sans Mono CJK KR}

\usepackage{amsmath,amssymb}
\usepackage{booktabs,longtable,tabularx,array}
\usepackage{enumitem}
\usepackage{xcolor}
\usepackage{hyperref}
\usepackage{fancyhdr}
\usepackage{microtype}
\usepackage[most]{tcolorbox}

\definecolor{CoreBlue}{HTML}{1F4E79}
\definecolor{CoreLight}{HTML}{EAF3F8}
\definecolor{StopRed}{HTML}{9C0006}
\definecolor{StopLight}{HTML}{FCE8E6}
\definecolor{GateGreen}{HTML}{38761D}
\definecolor{GateLight}{HTML}{EAF4E2}
\definecolor{DecisionOrange}{HTML}{B45F06}
\definecolor{DecisionLight}{HTML}{FFF2CC}

\hypersetup{
  colorlinks=true,
  linkcolor=CoreBlue,
  urlcolor=CoreBlue
}

\setlength{\parindent}{0pt}
\setlength{\parskip}{0.3em}
\setlength{\emergencystretch}{3em}
\setlist[itemize]{leftmargin=1.45em,itemsep=0.1em,topsep=0.2em}
\setlist[enumerate]{leftmargin=1.75em,itemsep=0.14em,topsep=0.22em}
\renewcommand{\arraystretch}{1.15}
\newcolumntype{L}[1]{>{\raggedright\arraybackslash}p{#1}}
\newcommand{\code}[1]{{\small\nolinkurl{#1}}}
\newcommand{\must}[1]{\textcolor{CoreBlue}{\textbf{#1}}}
\newcommand{\haltitem}[1]{\textcolor{StopRed}{\textbf{#1}}}
\newcommand{\passitem}[1]{\textcolor{GateGreen}{\textbf{#1}}}
\newcommand{\decisionitem}[1]{\textcolor{DecisionOrange}{\textbf{#1}}}

\newenvironment{contractbox}[1]
  {\begin{tcolorbox}[colback=CoreLight,colframe=CoreBlue,title={#1},fonttitle=\bfseries,before upper=\raggedright]}
  {\end{tcolorbox}}
\newenvironment{decisionbox}[1]
  {\begin{tcolorbox}[colback=DecisionLight,colframe=DecisionOrange,title={#1},fonttitle=\bfseries,before upper=\raggedright]}
  {\end{tcolorbox}}
\newenvironment{stopbox}[1]
  {\begin{tcolorbox}[colback=StopLight,colframe=StopRed,title={#1},fonttitle=\bfseries,before upper=\raggedright]}
  {\end{tcolorbox}}


\pagestyle{fancy}
\fancyhf{}
\lhead{Wind3DGS 개발 문서}
\rhead{R3 Multi-resplat Measure}
\cfoot{\thepage}

\title{R3 개발 계약\\Multi-resplat Measure Consistency}
\author{Wind3DGS}
\date{2026-08-22}

\begin{document}
\maketitle

\begin{contractbox}{문서 역할}
이 문서는 R3 구현과 실험을 진행하면서 계속 갱신하는 실행 계약이다.
방법의 수학적 의미와 claim의 최종 권위는 canonical idea sketch가 가지며,
본 문서는 R3에 필요한 입력, 구현 순서, 산출물과 중단 조건만 소유한다.
R0--R2 계약을 이 문서에서 임의로 재정의하지 않는다.
Canonical \code{eq:rd-area-mass}, \code{eq:rd-common-probe-map}과
\code{eq:rd-common-probe-adjoint}를 이름으로 참조한다.
\end{contractbox}

\DevChecklistPage{R3}{R2 선행 조건 대기}{
\DevProgressRow{0}{\hyperref[sec:r3-chapter-1]{1. 목적과 소유권}}{대기}{R2 이후 measure/latent 비교 범위와 owner 확인}
\DevProgressRow{0}{\hyperref[sec:r3-chapter-2]{2. 진입 입력}}{대기}{Independent GS variant와 동일 Teacher/probe 확보}
\DevProgressRow{0}{\hyperref[sec:r3-chapter-3]{3. 이미 동결된 계약}}{대기}{Map-law·공통 mask·measure·누설 방지 구현 확인}
\DevProgressRow{0}{\hyperref[sec:r3-chapter-4]{4. 구현 체크리스트}}{대기}{Pairing/measure/response consistency와 ablation 실행}
\DevProgressRow{0}{\hyperref[sec:r3-chapter-5]{5. Open Design Decisions}}{대기}{Density family, capacity와 Gate threshold 동결}
\DevProgressRow{0}{\hyperref[sec:r3-chapter-6]{6. 산출물}}{대기}{Variant registry, paired response와 Gate report 생성}
\DevProgressRow{0}{\hyperref[sec:r3-chapter-7]{7. 필수 fixture와 metric}}{대기}{Mapping floor·면적/질량·response/latent 비교 검증}
\DevProgressRow{0}{\hyperref[sec:r3-chapter-8]{8. 종료 조건과 실패 경로}}{대기}{Gate A\(_M\)/A\(_L\)과 direct-set fallback 경로 판정}
}{
\begin{contractbox}{진입 경계}Independent reconstruction과 GS별 density 비교는 아직 실행하지 않았다. 기존 개발용 Teacher probe 검사로 Gate A\(_M\)/A\(_L\)을 통과시킬 수 없다.\end{contractbox}
}

\tableofcontents

\section{목적과 소유권}
\label{sec:r3-chapter-1}

R3의 목적은 같은 source surface에서 독립적으로 복원한 3DGS들이 Gaussian 수, density,
seed와 view 조건이 달라도 \textbf{같은 물리적 표면 measure와 force-to-motion response}를
표현하는지 검증하는 것이다. R3는 다음 항목을 소유한다.

\begin{itemize}
  \item independent multi-resplat group의 생성 및 provenance 검증
  \item 동일한 canonical common-probe denominator에서의 paired 비교
  \item spatial surface measure, analytic traction integral, total force/work와 response dispersion 측정
  \item equal-weight, opacity-as-area와 consistency-loss 제거 ablation
  \item Gate A\(_M\) measure claim과 Gate A\(_L\) latent-representation claim의 증거
\end{itemize}

R3는 teacher convergence, force sign, exact-ZOH evaluator 또는 single-case Global overfit을
다시 승인하는 단계가 아니다. 이들은 R1--R2의 통과 결과를 입력으로 소비한다.

\begin{contractbox}{Primary evidence와 construction sanity의 구분}
\(\sum_i a_i=A_{\mathrm{ref}}\)와 \(\sum_i m_i=M_{\mathrm{ref}}\)는
정규화 식으로 보장되는 \textbf{construction sanity}다. 이 두 합이 맞는다는 사실만으로
density invariance나 novelty를 주장하지 않는다. R3의 primary evidence는 공간적으로 분해한
spatial surface measure, analytic traction 적분, force/work와 common-probe response의 independent-resplat
dispersion이다.
\end{contractbox}

\section{진입 입력}
\label{sec:r3-chapter-2}

R3 실행은 아래 입력이 모두 versioned artifact로 존재할 때만 시작한다.

\begin{longtable}{L{0.27\linewidth}L{0.61\linewidth}}
\toprule
입력 & 요구 내용\\
\midrule
\endfirsthead
\toprule
입력 & 요구 내용\\
\midrule
\endhead
\code{DomainContract} & SI 단위, champion domain, 허용 geometry/material/attachment/wind와 OOD 규칙\\
\code{GateATPrerequisiteReport} & R1 teacher/map/oracle과 R2 Global prerequisite가 통과한 identity/hash\\
\code{TeacherPackage} & 별도 spatial/temporal convergence를 통과한 teacher rest surface, quadrature와 trajectories\\
\code{CanonicalProbeMap} & teacher와 모든 declared GS variant가 공유하는 probe, ordered map와 하나의 frozen common-valid mask\\
\code{StaticGSPackage} & static GS payload, source-object ID, reconstruction provenance와 metric scale\\
\code{ResponsePackage} & R2 fixture를 통과한 area, derived mass, Global field/pole와 immutable evaluator identity\\
\code{R2GlobalCheckpoint} & R2 single-case 검증을 통과한 학습 코드와 초기 checkpoint; R3에서 사용한 정확한 hash 저장\\
\code{MetricRegistry} & metric ID, 단위, normalization, aggregation, 방향과 mapping-noise floor\\
\code{ThresholdRegistry} & Gate A\(_M\)/A\(_L\)의 dev-frozen tolerance와 Boolean 판정식\\
\bottomrule
\end{longtable}

같은 source object의 mesh, 모든 GS variant, material, wind sequence는 하나의 split group에 둔다.
Accepted Teacher의 시간·공간 통과는 동일 structure/element/mass/damping/attachment/load identity를
공유해야 한다. Consumer는 Teacher/GS 각각의 map-law/version을 확인하며 고차 signed Teacher map을
GS convex map으로 해석하지 않는다. R1의 개발용 sample은 이 accepted 입력을 대신하지 않는다.
Source surface 하나에서 한 GS를 단순 split/duplicate한 결과는 independent reconstruction으로 세지 않는다.

\section{이미 동결된 계약}
\label{sec:r3-chapter-3}

\subsection{Pair와 denominator}

\begin{itemize}
  \item Source object마다 2--3개의 predeclared density/seed/view/densification variant를 사용한다.
  정확한 nominal density ratio는 R3 Open Design Decision에서 run 전에 동결하며,
  실제 달성 density와 Gaussian 수는 nominal label과 분리해 기록한다.
  \item 모든 declared variant는 teacher와 같은 canonical probe set으로 remap한다.
  \item 정량 결과는 teacher와 group의 모든 declared variant가 함께 support하는
  \textbf{하나의 pre-test frozen common-valid mask}를 사용한다.
  \item 실패한 variant에 유리하도록 probe를 사후 삭제하거나 variant별 denominator를 만들지 않는다.
  \item constant/coordinate-linear reproduction, coverage와 mapping-noise floor를 통과하지 못한 map은
  predeclared reject/failure 규칙으로 처리한다.
\end{itemize}

\subsection{Measure와 force ownership}

Homogeneous V0는 positive area를 normalize하고 mass를 area에서 파생한다.
\[
  a_i=A_{\mathrm{ref}}
  \frac{\operatorname{softplus}(\xi_i^a)+\epsilon_a}
       {\sum_j\{\operatorname{softplus}(\xi_j^a)+\epsilon_a\}},
  \qquad
  m_i=\frac{M_{\mathrm{ref}}}{A_{\mathrm{ref}}}a_i.
\]
\(M_{\mathrm{ref}}\)만 total mass를 소유한다. Opacity, Gaussian count, teacher quadrature를
target package의 area/mass로 직접 대입하거나 asset별 scale correction을 하지 않는다.
R0에서 동결한 physical-valid policy를 모든 variant에 동일하게 적용한다.

Analytic traction은 area ownership 검사용 payload다. Probe traction \(\tau_P\)를 source total force로
보낼 때만 \(F_X=S_X^\top W_{A,P}\tau_P\)를 사용하며, total force payload에는
\(F_X=S_X^\top F_P\)를 사용한다. Runtime GS wind는 common-probe adjoint를 거치지 않고
각 Gaussian에서 \(F_i=a_i\tau_i\)로 직접 계산한다.

\subsection{누설 방지와 실험 동결}

\begin{itemize}
  \item Teacher--GS correspondence는 training loss와 evaluation fixture에만 사용한다.
  Target inference input이나 runtime package에 topology를 포함하지 않는다.
  \item Reconstruction 결과를 본 뒤 eligibility, density bin, support mask 또는 metric을 바꾸지 않는다.
  \item Probe peak/duration, analytic field family, wind program과 ablation budget은 development에서
  domain-wide로 동결한다.
  \item Report는 object-outer, variant-within-object paired 구조를 보존한다.
  Gaussian variant를 독립 object 표본처럼 세어 표본 수를 부풀리지 않는다.
\end{itemize}

\section{구현 체크리스트}
\label{sec:r3-chapter-4}

\subsection{Variant registry와 pairing}

\begin{itemize}
  \item[$\square$] Source-object ID와 split-group ID를 모든 teacher/GS artifact에 기록
  \item[$\square$] Variant별 reconstruction provenance 저장: camera subset, seed, initialization과 source commit
  \item[$\square$] Densification schedule, nominal/achieved density와 achieved Gaussian count를 별도 field로 저장
  \item[$\square$] Independent reconstruction 판정기를 만들고 split/duplicate 파생본을 거부
  \item[$\square$] Declared variant 목록을 map/metric 실행 전에 freeze하고 group hash 생성
  \item[$\square$] Static-GS eligibility 결과와 제외 사유를 denominator와 함께 저장
\end{itemize}

\subsection{Spatial surface measure와 analytic quadrature}

\begin{itemize}
  \item[$\square$] Frozen canonical probe partition 또는 predeclared surface cells에서
  predicted local area와 derived mass를 aggregate
  \item[$\square$] Constant 및 coordinate-linear scalar measure field의 local integral 오차 측정
  \item[$\square$] Constant, linear-gradient와 piecewise-localized analytic traction family를
  모든 variant에 동일하게 적용
  \item[$\square$] Net force, first moment(정의된 경우), virtual work와 teacher quadrature 대비 오차 측정
  \item[$\square$] Total area/mass equality는 sanity column으로 분리하고 primary score에 중복 가산하지 않음
  \item[$\square$] Physical-valid mask의 면적 누락, background/floater load와 zero-row 위반 검사
\end{itemize}

\subsection{Response consistency와 ablation}

\begin{itemize}
  \item[$\square$] Impulse, step-on/off, multi-sine/log-chirp와 traveling gust를 동일 physical-time grid에서 실행
  \item[$\square$] Common-probe displacement/velocity, tip, amplitude, phase, PSD와 external work 저장
  \item[$\square$] Pairwise뿐 아니라 group centroid 대비 dispersion과 worst-variant degradation 보고
  \item[$\square$] Mapping-noise floor 이하/초과 성분을 분리해 model error로 과대 해석하지 않음
  \item[$\square$] Equal Gaussian weight, opacity-as-area, consistency loss 제거,
  latent/relation auxiliary 제거, patch-only ablation 실행
  \item[$\square$] Teacher quadrature를 직접 쓰는 inference-impossible oracle을 ceiling으로 명시
  \item[$\square$] Gaussian 수 대비 setup memory/latency와 evaluator memory slope 기록
\end{itemize}

\subsection{Latent와 direct-set route 판정}

\begin{itemize}
  \item[$\square$] Latent-relation model과 capacity-matched direct-set response model을 동일 split/probe/evaluator로 학습
  \item[$\square$] Parameter, setup FLOPs, memory, training budget와 dev checkpoint selection rule을 사전 동결
  \item[$\square$] Direct-set은 latent token/relation auxiliary 없이 permutation-equivariant GS set에서 package를 직접 예측
  \item[$\square$] Single-case R2 진단이 아니라 paired multi-resplat dev evidence로 Gate A\(_L\) 판정
  \item[$\square$] A\(_L\) pass면 latent route, fail이면 direct-set route의 config/checkpoint를 R4 hand-off로 발행
  \item[$\square$] 선택 route의 encoder family와 capacity-matching 조건을 보존하는 유한 rank/pole/capacity
  후보를 \code{GlobalConfigSet}으로 hash; R2b loss/schedule은 후보에서 제외
\end{itemize}

\subsection{Gate 계산}

\begin{itemize}
  \item[$\square$] Metric별 단위, normalization, object/variant/time reduction과 pass 방향을 registry에서 로드
  \item[$\square$] Gate A\(_M\)와 A\(_L\) Boolean 식을 코드에 하드코딩하지 않고
  frozen registry hash로 평가
  \item[$\square$] A\(_T\) prerequisite hash를 검증하고 실패 시 A\(_M\)/A\(_L\) 수치로 pipeline을 열지 않음
  \item[$\square$] Mean뿐 아니라 worst declared density/seed variant와 method success rate 보고
  \item[$\square$] Exclusion count와 이유를 success denominator에 포함
  \item[$\square$] 결과를 본 뒤 threshold, primary metric 또는 density bin 변경 금지
\end{itemize}

\section{Open Design Decisions}
\label{sec:r3-chapter-5}

\begin{decisionbox}{R3 시작 전에 닫을 결정}
다음 항목은 canonical 방법 identity가 아니라 R3 개발 계약의 미결정값이다.
선택 근거, freeze 시점과 config hash를 decision log에 남긴다.
\end{decisionbox}

\begin{enumerate}
  \item \decisionitem{Static-GS reconstruction eligibility:}
  최소 coverage, opacity/floater 품질, metric-scale 신뢰도와 map 지원률을 어떤 Boolean 규칙으로 묶을지 결정한다.
  Eligibility는 teacher response 또는 test motion error를 보지 않아야 한다.
  \item \decisionitem{Achieved count/density registry:}
  nominal target density와 실제 Gaussian count, surface coverage 및 local sampling density를 어떤 schema로
  기록하고 density bin을 어떻게 확정할지 결정한다. \(0.5\times,1\times,2\times\)는 권장 시작점일 뿐
  동결값이 아니며 reconstruction feasibility를 본 development rule로 run 전에 확정한다.
  \item \decisionitem{최소 variant 수:}
  source object당 density/seed 조합과 reconstruction 재시도 한도를 정한다.
  실패한 reconstruction을 성공할 때까지 무제한 교체하지 않는다.
  \item \decisionitem{Local measure partition:}
  canonical probe patch, fixed UV/arc-length bin 또는 geometric cell 중 어떤 partition을 primary로 사용할지 결정한다.
  \item \decisionitem{Analytic traction family:}
  constant/linear/localized field의 정확한 식, 크기, 방향과 normalization을 정한다.
  \item \decisionitem{Gate A 수치:}
  primary metric, object-outer 집계, worst-case allowance와 A\(_M\)/A\(_L\) Boolean threshold를 dev에서 동결한다.
\end{enumerate}

\section{산출물}
\label{sec:r3-chapter-6}

\begin{longtable}{L{0.30\linewidth}L{0.58\linewidth}}
\toprule
산출물 & 필수 내용\\
\midrule
\endfirsthead
\toprule
산출물 & 필수 내용\\
\midrule
\endhead
\code{MultiResplatGroupManifest} & source/split ID, declared variants, independence provenance, group hash\\
\code{StaticGSVariantRecord} & nominal/achieved density, Gaussian count, eligibility, map coverage와 exclusion reason\\
\code{SpatialMeasureReport} & spatial area/mass, analytic traction integral, net force/first moment/work 오차\\
\code{ResponseDispersionReport} & common-probe paired/group dispersion, worst variant, mapping-noise 분해\\
\code{R3AblationReport} & full/equal/opacity/no-consistency/direct-set/patch-only/oracle 비교\\
\code{GateAReport} & A\(_T\) prerequisite hash, A\(_M\)/A\(_L\) Boolean 입력과 최종 판정\\
\code{GlobalConfigSet} & 선택 route 안의 유한 Global rank/pole/capacity/compute 후보와 capacity-matching invalidation hash\\
\code{R3RouteConfig} & latent 또는 direct-set main route, fixed-family 제한, architecture/config와 invalidation hash\\
\code{R3RouteCheckpoint} & 선택 route의 dev-selected checkpoint와 R4 initialization/training provenance\\
\bottomrule
\end{longtable}

모든 report는 source commit과 model, package, config identity를 포함한다.
Dataset split, metric/threshold registry, probe/map/common-valid mask와 evaluator hash도 함께 기록한다.

\section{필수 fixture와 metric}
\label{sec:r3-chapter-7}

\subsection{필수 fixture}

\begin{enumerate}
  \item \textbf{Construction sanity:} total area/mass, nonnegative area/mass, invalid row와 unit 검사
  \item \textbf{Spatial measure reproduction:} constant/coordinate-linear field의 patch별 적분
  \item \textbf{Analytic traction:} spatially constant/linear/localized traction의 force, moment와 work
  \item \textbf{Paired response:} impulse, step/recovery, frequency probe와 traveling gust
  \item \textbf{Ablation parity:} 동일 split, budget, probe mask와 evaluator로 모든 measure baseline 비교
  \item \textbf{Scaling:} achieved Gaussian count 대비 setup/package/runtime memory와 latency
\end{enumerate}

\subsection{Primary metric}

\begin{itemize}
  \item spatial area/mass normalized error와 independent-variant dispersion
  \item analytic traction의 net-force/virtual-work normalized error
  \item external work dispersion과 common-probe velocity/response dispersion
  \item density별 phase/amplitude/PSD error와 worst-density degradation
\end{itemize}

Total area/mass equality, positivity와 hash 일치는 필수 sanity이지만 primary novelty metric이 아니다.
Token coverage/duplicate/collapse, map coverage와 latency는 mandatory diagnostic이다.

\section{종료 조건과 실패 경로}
\label{sec:r3-chapter-8}

\begin{contractbox}{R3 종료 조건}
Gate A\(_T\) prerequisite를 통과한 뒤 declared multi-resplat group 전체가 frozen denominator로 평가되고,
spatial surface measure, analytic traction, force/work와 response dispersion report가 재현 가능하며,
Gate A\(_M\)/A\(_L\) Boolean 판정이 frozen registry로 생성되면 R3를 종료한다.
A\(_M\)/A\(_L\) 결과에 맞는 \code{R3RouteConfig}/\code{R3RouteCheckpoint}를 발행하면 R4로 진행할 수 있으며,
density-aware held-out claim은 A\(_M\) pass일 때만 허용한다.
\end{contractbox}

\begin{itemize}
  \item \passitem{A\(_M\) pass, A\(_L\) pass:} measure-consistent latent response claim을 유지하고 R4로 진행한다.
  \item \haltitem{A\(_M\) fail only:} density-invariant claim을 철회하고 fixed reconstruction family로 축소한다.
  R4 진행 여부와 허용 reconstruction family를 새 hash로 고정한다.
  \item \haltitem{A\(_L\) fail only:} latent-anchor/relation contribution을 삭제하고 capacity-matched direct-set
  response model을 R4의 main route로 사용한다.
  \item \haltitem{A\(_M\), A\(_L\) both fail:} fixed reconstruction family로 제한한 direct-set response를
  main route로 발행하고 density/latent contribution을 모두 철회한다.
  \item \haltitem{A\(_T\) prerequisite fail:} R3 수치로 덮지 않고 R1--R2로 되돌아간다.
  \item \haltitem{Eligibility collapse:} declared density/seed group 대부분이 static-GS eligibility를 통과하지 못하면
  denominator를 사후 축소하지 않고 reconstruction domain 또는 방법 claim을 재설계한다.
  \item \haltitem{Mapping-floor domination:} primary response 차이가 mapping-noise floor로 설명되면
  density claim의 evidence로 사용하지 않고 map/probe resolution을 먼저 재검토한다.
\end{itemize}

\end{document}
